\documentclass[10pt,a4paper]{amsart}
\usepackage[margin=19mm]{geometry}
\usepackage{amsmath,amssymb,mathtools}
\usepackage[scr=rsfs,cal=euler]{mathalfa}
\usepackage{xfrac}
\usepackage[unicode,hidelinks,bookmarks=false]{hyperref}
\newcommand{\ie}{i.e.\ }
\newcommand{\cf}{cf.\ }
\newcommand{\resp}{respectively\ }
\newcommand{\Subsection}{\subsection}
\providecommand{\nobreakdash}{\nobreak\hbox{-}}
\setlength{\emergencystretch}{2em}
\multlinegap0pt
\usepackage{mathtools}
\usepackage{amssymb}
\usepackage{enumerate}
\usepackage[shortlabels]{enumitem}
\usepackage{xcolor}
\usepackage{algorithm}\usepackage{algpseudocode}
\newtheorem{proposition}{Proposition}
\title{A discrete-time overdetermined problem for the heat equation}
\author{Lorenzo Cavallina, Andrea Pinamonti}
\date{Source: arXiv:2604.20430v1 (CC BY 4.0)}
\begin{document}
\maketitle

\noindent\emph{Source and licence.} A discrete-time overdetermined problem for the heat equation. Version arXiv:2604.20430v1 (CC BY 4.0), distributed by arXiv under the Creative Commons Attribution 4.0 International licence, http://creativecommons.org/licenses/by/4.0/, as recorded in the arXiv licence metadata for this exact version and shown on the abstract page https://arxiv.org/abs/2604.20430v1. Full licence notice: \texttt{licenses/arxiv-2604.20430-CC-BY-4.0.txt}.\par\smallskip

\noindent\textit{Editorial scope.} Counted unit: the article's proposition prop:riem-cfp (source0.tex lines 865--880). The article's complete original proof of it is delivered with it (source0.tex lines 882--938, one \texttt{proof} environment of 57 lines). No mathematics is written, repaired or re-derived: the statement and the proof are the article's own lines, in the article's own notation, and no retained line is substituted or re-flowed.  Retained uncounted as the source-local context the proof needs: lines 847--850.  Nothing was omitted from any retained span, so the package carries no omission record.  No retained line contains a citation, so the wrapper carries no bibliography and no bibliography entry.  No retained line contains a figure environment, an image-inclusion command or drawing code, so no image is delivered and the package declares no asset dependency.  Editorial formatting only, in addition: an AMS \texttt{amsart} preamble that restates the article's own declarations and macros used by the retained lines, namely the environments \texttt{proposition} on separate counters, together with the standard packages those lines need.  The retained environments print in this document as Proposition 1 in order of appearance, whereas the article prints the same items with the numbering of its own full document.  The complete native source of the article as distributed in the arXiv e-print for this exact version is bundled unchanged as \texttt{sources/198933/source0.tex}.
\begin{equation}\label{eq:spectral-du}
\partial_\nu u(t)=\sum_{j=1}^\infty a_j e^{-\lambda_j t}\partial_\nu\phi_j
\qquad\text{in }C^\infty(\partial\Omega)\text{ for }t\geq \tau,
\end{equation}

\begin{proposition}\label{prop:riem-cfp}
Assume that there exists a sequence $(t_m)_m\subset (0,\infty)$ such that
\[
\partial_\nu u(t_m)\in \mathcal C(\partial\Omega)\qquad\text{for all }m\in\mathbb{N},
\]
and either
\[
t_m\to t_\ast\in (0,\infty),
\qquad\text{or}\qquad
t_m\to \infty.
\]
Then $\Omega$ has the \emph{constant flow property}, namely
\[
\partial_\nu u(t)\in \mathcal C(\partial\Omega)\qquad\text{for every fixed }t>0.
\]
\end{proposition}

\begin{proof}
Fix $\psi\in C^\infty(\partial\Omega)$ with zero average on $\partial\Omega$ and
set
\[
F_\psi(t):=\int_{\partial\Omega}\psi\,\partial_\nu u(t).
\]
By \eqref{eq:spectral-du},
\begin{equation}\label{eq:Fpsi-series}
F_\psi(t)=\sum_{j=1}^\infty b_j e^{-\lambda_j t},
\qquad
b_j:=a_j\int_{\partial\Omega}\psi\,\partial_\nu\phi_j,
\end{equation}
with absolute and uniform convergence on $[\tau,\infty)$ for every $\tau>0$.
Hence $F_\psi$ is real-analytic on $(0,\infty)$.

If $t_m\to t_\ast\in(0,\infty)$, then $F_\psi(t_m)=0$ for every $m\in\mathbb{N}$ by the
assumption $\partial_\nu u(t_m)\in\mathcal C(\partial\Omega)$. Since the
zeros accumulate at an interior point of $(0,\infty)$ and $F_\psi$ is
real-analytic, we conclude that $F_\psi\equiv 0$ on $(0,\infty)$.

Assume now that $t_m\to\infty$. Let
\[
0<\Lambda_1<\Lambda_2<\Lambda_3<\cdots
\]
be the distinct Dirichlet eigenvalues of $\Delta_g$ on $\Omega$, and rewrite
\eqref{eq:Fpsi-series} as
\[
F_\psi(t)=\sum_{\ell=1}^\infty \beta_\ell e^{-\Lambda_\ell t},
\qquad
\beta_\ell:=\sum_{\lambda_j=\Lambda_\ell} b_j.
\]
We claim that $\beta_\ell=0$ for every $\ell$. Since the series converges
absolutely on $[T,\infty)$ for every $T>0$, we can choose $T>0$ so that
$t_m\geq 2T$ for all $m$ large enough. Then
\[
0=e^{\Lambda_1 t_m}F_\psi(t_m)
=\beta_1+\sum_{\ell\geq 2}\beta_\ell e^{-(\Lambda_\ell-\Lambda_1)t_m}.
\]
Moreover,
\begin{align*}
\sum_{\ell\geq 2} |\beta_\ell| e^{-(\Lambda_\ell-\Lambda_1)t_m}
&=
\sum_{\ell\geq 2} |\beta_\ell| e^{-\Lambda_\ell T}
e^{-\Lambda_\ell(t_m-T)}e^{\Lambda_1 t_m}\\
&\leq
e^{-(\Lambda_2-\Lambda_1)t_m+\Lambda_2 T}
\sum_{\ell\geq 2} |\beta_\ell|e^{-\Lambda_\ell T}\xrightarrow[m\to\infty]{}0.
\end{align*}
Hence $\beta_1=0$. Subtracting the first term and iterating the same argument,
we obtain $\beta_\ell=0$ for every $\ell$, hence $F_\psi\equiv 0$ also in this
case.

Since $F_\psi(t)=0$ for every zero-average $\psi$ and every $t>0$, it follows
that $\partial_\nu u(t)$ is orthogonal to all zero-average smooth
functions on $\partial\Omega$, and therefore
$\partial_\nu u(t)\in\mathcal C(\partial\Omega)$ for every $t>0$.
\end{proof}

\end{document}
